\documentclass[a4paper,11pt]{article}
\usepackage[utf8]{inputenc}
\usepackage[T1]{fontenc}
\usepackage[ngerman]{babel}
\usepackage{geometry}
\geometry{left=2.5cm, right=2.5cm, top=2.5cm, bottom=2.5cm}
\usepackage{longtable}
\usepackage{booktabs}
\usepackage{array}
\usepackage{xcolor}
\usepackage{colortbl}
\usepackage{hyperref}
\usepackage{fancyhdr}
\usepackage{tabularx}
\usepackage{enumitem}
\usepackage{graphicx}
\usepackage{listings}

\definecolor{headerblue}{HTML}{1565C0}
\definecolor{tableheader}{HTML}{E3F2FD}
\definecolor{pass}{HTML}{2E7D32}
\definecolor{warn}{HTML}{F57F17}
\definecolor{codebg}{HTML}{F5F5F5}

\hypersetup{
    colorlinks=true,
    linkcolor=headerblue,
    urlcolor=headerblue
}

\pagestyle{fancy}
\fancyhf{}
\fancyhead[L]{\textbf{INSPECTAIR}}
\fancyhead[R]{Configuration Management v2.0}
\fancyfoot[C]{\thepage}
\renewcommand{\headrulewidth}{0.4pt}

\setlength{\parindent}{0pt}
\setlength{\parskip}{6pt}

\lstset{
    backgroundcolor=\color{codebg},
    basicstyle=\ttfamily\small,
    breaklines=true,
    frame=single,
    rulecolor=\color{gray!40},
    xleftmargin=2em,
    framexleftmargin=1.5em,
    showstringspaces=false,
    tabsize=4
}

\begin{document}

\begin{center}
{\LARGE\textbf{INSPECTAIR}}\\[0.3cm]
{\Large Configuration Management}\\[0.3cm]
{\large Luftqualitätsmessgerät mit ESP32-S3}\\[0.5cm]
\begin{tabular}{ll}
\textbf{Projekt:} & InspectAir -- Luftqualitätsmessgerät \\
\textbf{Version:} & 2.0 \\
\textbf{Datum:} & Februar 2026 \\
\textbf{Autoren:} & Maximilian Liebl, Cedric Stroh, Jannis Messer \\
\textbf{Modul:} & Mikrocomputertechnik 1, DHBW Ravensburg \\
\textbf{Dozent:} & Thomas Stingl (Airbus) \\
\end{tabular}
\end{center}

\vspace{0.5cm}
\noindent\textit{Dieses Dokument wurde unter Verwendung von KI-Tools (Claude, Anthropic) zur Überprüfung erstellt.}
\vspace{0.3cm}

\tableofcontents
\newpage

% ============================================================
\section{Versionskontrolle}
% ============================================================

\begin{longtable}{p{4cm} p{9cm}}
\toprule
\rowcolor{tableheader}
\textbf{Eigenschaft} & \textbf{Details} \\
\midrule
\endhead
Tool       & Git \\
Repository & GitHub (privat) \\
Branching  & \texttt{main} + Feature-Branches \\
\bottomrule
\end{longtable}

% ============================================================
\section{Entwicklungsumgebung}
% ============================================================

\begin{longtable}{p{4cm} p{9cm}}
\toprule
\rowcolor{tableheader}
\textbf{Tool} & \textbf{Version / Details} \\
\midrule
\endhead
IDE            & Visual Studio Code \\
Build-System   & PlatformIO \\
Framework      & Arduino (espressif32 Platform) \\
Board          & esp32-s3-devkitc-1 \\
Upload Speed   & 921600 Baud \\
Monitor Speed  & 115200 Baud \\
\bottomrule
\end{longtable}

% ============================================================
\section{platformio.ini}
% ============================================================

Zentrale Build-Konfiguration des Projekts:

\begin{lstlisting}
[env:esp32-s3-devkitc-1]
platform = espressif32
board = esp32-s3-devkitc-1
framework = arduino
monitor_speed = 115200
monitor_rts = 0
monitor_dtr = 0
board_upload.flash_size = 8MB
board_build.arduino.memory_type = qio_opi
build_flags =
    -DARDUINO_USB_CDC_ON_BOOT=1
    -DARDUINO_USB_MODE=1
    -DBOARD_HAS_PSRAM
    -DLV_CONF_SKIP=1
    -DLV_COLOR_DEPTH=16
    -DLV_COLOR_16_SWAP=1
    -DLV_MEM_SIZE="(96*1024)"
    ...
lib_deps =
    Adafruit_AHTX0 (GitHub)
    Adafruit_SGP40 (GitHub)
    fu-hsi/PMS Library @ ^1.1.0
    wifwaf/MH-Z19 @ ^1.5.4
    ncmreynolds/ld2410 @ ^0.1.3
    lovyan03/LovyanGFX @ ^1.1.16
    lvgl/lvgl @ ^9.2.2
    plerup/EspSoftwareSerial @ ^8.2.0
\end{lstlisting}

\subsection{Bibliotheken-Übersicht}

\begin{longtable}{p{4.5cm} p{3cm} p{5.5cm}}
\toprule
\rowcolor{tableheader}
\textbf{Bibliothek} & \textbf{Version} & \textbf{Verwendung} \\
\midrule
\endhead
LovyanGFX       & 1.1.16 & Display-Treiber (ST7796S via SPI) \\
LVGL             & 9.2.2  & UI-Framework (5 Screens) \\
Adafruit AHTX0   & GitHub & AHT20 Temperatur/Feuchte \\
Adafruit SGP40   & GitHub & SGP40 VOC-Sensor \\
PMS Library      & 1.1.0  & PMS5003 Feinstaub \\
MH-Z19           & 1.5.4  & MH-Z19C CO\textsubscript{2}-Sensor \\
ld2410           & 0.1.3  & LD2410C Radar \\
EspSoftwareSerial & 8.2.0 & SoftwareSerial für CO\textsubscript{2} \\
\bottomrule
\end{longtable}

% ============================================================
\section{Hardware-Konfiguration}
% ============================================================

Alle Pin-Definitionen sind zentral in \texttt{pins.h} definiert.

\subsection{Display ST7796S (SPI)}

\begin{longtable}{p{3cm} p{2.5cm} p{7.5cm}}
\toprule
\rowcolor{tableheader}
\textbf{Signal} & \textbf{GPIO} & \textbf{Hinweis} \\
\midrule
\endhead
TFT\_MOSI & 11 & SPI Data \\
TFT\_MISO & 13 & SPI Data (nicht aktiv genutzt) \\
TFT\_SCLK & 12 & SPI Clock \\
TFT\_CS   & 14 & Chip Select \\
TFT\_DC   & 9  & Data/Command \\
TFT\_RST  & 46 & Reset \\
TFT\_BL   & 3  & Backlight (via BS170 MOSFET) \\
\bottomrule
\end{longtable}

\subsection{I2C Sensoren (AHT20, SGP40)}

\begin{longtable}{p{3cm} p{2.5cm} p{7.5cm}}
\toprule
\rowcolor{tableheader}
\textbf{Signal} & \textbf{GPIO} & \textbf{Hinweis} \\
\midrule
\endhead
I2C\_SDA & 8  & Data \\
I2C\_SCL & 18 & Clock \\
\bottomrule
\end{longtable}

\subsection{UART Sensoren}

\begin{longtable}{p{3cm} p{2cm} p{2cm} p{5.5cm}}
\toprule
\rowcolor{tableheader}
\textbf{Sensor} & \textbf{RX} & \textbf{TX} & \textbf{Hinweis} \\
\midrule
\endhead
PMS5003  & 16 & 17 & UART1, 9600 Baud \\
MH-Z19C  & 4  & 5  & SoftwareSerial, 9600 Baud \\
LD2410C  & 7  & 6  & UART2, 256000 Baud, Level Shifter! \\
\bottomrule
\end{longtable}

\textbf{Hinweis:} RX/TX beziehen sich auf die ESP32-Perspektive (RX = ESP empfängt, TX = ESP sendet).

\subsection{Sonstige Pins}

\begin{longtable}{p{3.5cm} p{2.5cm} p{7cm}}
\toprule
\rowcolor{tableheader}
\textbf{Funktion} & \textbf{GPIO} & \textbf{Hinweis} \\
\midrule
\endhead
RADAR\_OUT & 15 & LD2410C Presence Output (Digital) \\
UI\_BUTTON & 1  & Screen-Wechsel (Active LOW, interner Pullup) \\
\bottomrule
\end{longtable}

% ============================================================
\section{Dokumenten-Versionen}
% ============================================================

\begin{longtable}{p{5cm} p{2cm} p{3cm} p{3cm}}
\toprule
\rowcolor{tableheader}
\textbf{Dokument} & \textbf{Version} & \textbf{Datum} & \textbf{Autor} \\
\midrule
\endhead
Requirements Specification   & 2.0 & Feb 2026 & Team \\
Design Description           & 2.0 & Feb 2026 & Team \\
Test Specification           & 2.0 & Feb 2026 & Maximilian Liebl \\
Schedule                     & 2.0 & Feb 2026 & Team \\
Risk \& Opportunities        & 2.0 & Feb 2026 & Team \\
Coding Rules                 & 2.0 & Feb 2026 & Maximilian Liebl \\
Configuration Management     & 2.0 & Feb 2026 & Maximilian Liebl \\
Pin-Belegung \& Bauteilliste & 2.0 & Feb 2026 & Team \\
\bottomrule
\end{longtable}

\end{document}
